\documentclass[11pt,a4paper,oneside]{report}

% ---------------------------------------------------------------------------
% Jazyk, kódování a písma
% ---------------------------------------------------------------------------
\usepackage[utf8]{inputenc}
\usepackage[T1]{fontenc}
\usepackage[czech]{babel}
\usepackage{lmodern}
\usepackage{microtype}

% ---------------------------------------------------------------------------
% Geometrie stránky a záhlaví
% ---------------------------------------------------------------------------
\usepackage[
    a4paper,
    top=2.5cm,
    bottom=2.5cm,
    left=2.5cm,
    right=2.5cm,
    headheight=14pt
]{geometry}

\usepackage{fancyhdr}
\pagestyle{fancy}
\fancyhf{}
\fancyhead[L]{\small\nouppercase{\leftmark}}
\fancyhead[R]{\small\thepage}
\renewcommand{\headrulewidth}{0.4pt}

% ---------------------------------------------------------------------------
% Matematická sazba a fyzikální symboly
% ---------------------------------------------------------------------------
\usepackage{amsmath,amssymb,amsfonts,amsthm}
\usepackage{bm}
\usepackage{siunitx}
\sisetup{output-decimal-marker = {,}}

% Pomocná makra pro matematiku a mechaniku
\providecommand{\coloneqq}{\mathrel{\mathop:}=}
\newcommand{\vect}[1]{\bm{#1}}                       % Vektor (tučně)
\newcommand{\dif}{\mathop{}\!\mathrm{d}}             % Diferenciál 'd'
\newcommand{\odv}[2]{\frac{\dif #1}{\dif #2}}        % Obyčejná derivace
\newcommand{\dodv}[2]{\frac{\dif^2 #1}{\dif #2^2}}   % Druhá obyčejná derivace
\newcommand{\pdv}[2]{\frac{\partial #1}{\partial #2}}% Parciální derivace
\newcommand{\dpdv}[2]{\frac{\partial^2 #1}{\partial #2^2}} % Druhá parciální derivace
\newcommand{\dotprod}{\bm{\cdot}}                    % Skalární součin
\newcommand{\crossprod}{\bm{\times}}                 % Vektorový součin
\newcommand{\Lagr}{\mathcal{L}}                      % Lagrangeián
\newcommand{\Ham}{\mathcal{H}}                       % Hamiltonián

% ---------------------------------------------------------------------------
% Barvy a grafika
% ---------------------------------------------------------------------------
\usepackage{graphicx}
\usepackage{booktabs}
\usepackage{xcolor}

\definecolor{primary}{RGB}{18, 52, 86}      % Tmavě modrá
\definecolor{secondary}{RGB}{140, 40, 40}   % Cihlově červená
\definecolor{boxbg}{RGB}{248, 249, 252}     % Světle šedé pozadí
\definecolor{boxborder}{RGB}{210, 215, 225} % Rámeček
\definecolor{notebg}{RGB}{255, 253, 240}    % Pozadí pro poznámku
\definecolor{noteborder}{RGB}{230, 200, 100}

% ---------------------------------------------------------------------------
% Rámečky pro věty, definice a příklady (mdframed)
% ---------------------------------------------------------------------------
\usepackage[framemethod=tikz]{mdframed}

\newtheoremstyle{mystyle}%
  {6pt}{6pt}%
  {\normalfont}%
  {}%
  {\bfseries\color{primary}}%
  {.}%
  {0.5em}%
  {\thmname{#1}\thmnumber{ #2}\thmnote{ (#3)}}

\theoremstyle{mystyle}
\newtheorem{definicethm}{Definice}[chapter]
\newtheorem{vetathm}{Věta}[chapter]
\newtheorem{prikladthm}{Příklad}[chapter]

\newenvironment{definice}[1][]{%
    \begin{mdframed}[
        roundcorner=4pt,
        backgroundcolor=boxbg,
        linecolor=primary,
        linewidth=1.2pt,
        innertopmargin=8pt,
        innerbottommargin=8pt,
        innerleftmargin=12pt,
        innerrightmargin=12pt,
        skipabove=10pt,
        skipbelow=10pt,
        nobreak=false
    ]%
    \begin{definicethm}[#1]%
}{%
    \end{definicethm}%
    \end{mdframed}%
}

\newenvironment{veta}[1][]{%
    \begin{mdframed}[
        roundcorner=4pt,
        backgroundcolor=boxbg,
        linecolor=secondary,
        linewidth=1.2pt,
        innertopmargin=8pt,
        innerbottommargin=8pt,
        innerleftmargin=12pt,
        innerrightmargin=12pt,
        skipabove=10pt,
        skipbelow=10pt,
        nobreak=false
    ]%
    \begin{vetathm}[#1]%
}{%
    \end{vetathm}%
    \end{mdframed}%
}

\newenvironment{priklad}[1][]{%
    \begin{mdframed}[
        topline=false,
        rightline=false,
        bottomline=false,
        linecolor=gray,
        linewidth=2.5pt,
        innerleftmargin=10pt,
        innerrightmargin=0pt,
        innertopmargin=4pt,
        innerbottommargin=4pt,
        skipabove=8pt,
        skipbelow=8pt
    ]%
    \begin{prikladthm}[#1]%
}{%
    \end{prikladthm}%
    \end{mdframed}%
}

\newenvironment{poznamka}{%
    \begin{mdframed}[
        roundcorner=3pt,
        backgroundcolor=notebg,
        linecolor=noteborder,
        linewidth=0.8pt,
        innertopmargin=6pt,
        innerbottommargin=6pt,
        innerleftmargin=10pt,
        innerrightmargin=10pt,
        skipabove=8pt,
        skipbelow=8pt
    ]%
    \small\textbf{\color{noteborder!70!black}Poznámka: }%
}{%
    \end{mdframed}%
}

% ---------------------------------------------------------------------------
% Odkazy a metadata PDF
% ---------------------------------------------------------------------------
\usepackage{hyperref}
\hypersetup{
    colorlinks=true,
    linkcolor=primary,
    citecolor=secondary,
    urlcolor=blue!70!black,
    pdftitle={Klasická mechanika -- Zápisky z přednášek},
    pdfauthor={Jméno Studenta / Přednášejícího},
    pdfsubject={Klasická teoretická mechanika}
}

\begin{document}

% ---------------------------------------------------------------------------
% Titulní strana
% ---------------------------------------------------------------------------
\pagenumbering{Alph}
\begin{titlepage}
    \centering
    \vspace*{2cm}
    {\Huge\bfseries Mechanika a kontinuum\par}
    \vspace{0.8cm}
    {\Large Zápisky z přednášek\par}
    \vspace{0.3cm}
    {\large NAFY001\par}
    \vfill
    {\textbf{Autor: } Emil Varga\par}
    \vspace{0.2cm}
    {\small Univerzita Karlova MFF\par}
    {\small \texttt{emil.varga@matfyz.cuni.cz}\par}
    \vspace{1.5cm}
    {\large Akademický rok 2026/2027\par}
\end{titlepage}

% ---------------------------------------------------------------------------
% Předmluva a obsah (římské číslování)
% ---------------------------------------------------------------------------
\pagenumbering{roman}

\tableofcontents
\newpage

% ---------------------------------------------------------------------------
% Úvod
% ---------------------------------------------------------------------------
\chapter*{Úvod}
\addcontentsline{toc}{chapter}{Úvod}

\section*{Vědecký obor}

Fyzika studuje obecné vlastnosti látek a~polí. Fyzika je \emph{vědecký obor}, a~proto všechna tvrzení v~ní musejí být přímo nebo nepřímo, alespoň v~principu, vyvratitelná experimentem (falzifikovatelná).

Konstrukce fyzikálních teorií se obecně řídí tzv.~\emph{Occamovou břitvou}: modely a tvrzení by měly být co nejjednodušší, avšak v~souladu se všemi dostupnými pozorováními. Není-li nějaká část teorie nezbytná pro vysvětlení pozorovaných jevů, z~teorie se odstraní.

Fyzika a chemie mají široký překryv. Atomy a molekuly se řídí zákony kvantové mechaniky, i~když odvození zákonů chemie přímo ze základních zákonů fyziky představuje silně netriviální problém mnoha částic.

\subsection*{Mechanika}

Hrubé, nejednoznačné a nepřesné, zato však běžně používané dělení fyzikálních oborů je například:
\begin{description}
    \item[mechanika] pohyb těles a jejich vzájemné silové působení,
    \item[termodynamika a statistická fyzika] jevy způsobené chaotickým pohybem velkého množství částic,
    \item[elektřina a magnetismus] elektromagnetické pole a jeho interakce s~hmotou,
    \item[jaderná a částicová fyzika] jevy v~atomovém jádře a struktura mikrosvěta.
\end{description}

Mechanika je obor fyziky, který se zabývá pohybem těles a jejich vzájemným působením. Samotným popisem pohybu bez ohledu na jeho příčiny se zabývá \emph{kinematika}. Příčiny pohybu, tedy proč pohyb vzniká nebo se mění v~čase vlivem sil, jsou předmětem \emph{dynamiky}.

V~prvním přiblížení se budeme věnovat pohybu \emph{hmotných bodů}, tj.~idealizaci reálných objektů, které lze plně popsat pouze jejich hmotností, polohou a rychlostí. Hmotné body nemají žádnou vnitřní strukturu ani vlastní rozměry. \emph{Klasická (newtonovská)} mechanika popisuje pohyb těles, která jsou dostatečně těžká ve srovnání s~atomy, ale ne natolik hmotná, aby bylo nutné uvažovat zakřivení prostoročasu, a která se na makroskopických škálách pohybují rychlostmi podstatně menšími, než je rychlost světla.

Newtonovská mechanika selhává, když se rychlosti těles začnou blížit rychlosti světla (čímž se dostáváme do \emph{speciální teorie relativity}) nebo pokud jsou gravitační pole příliš silná (což je doména \emph{obecné teorie relativity}).

Naopak při malých hmotnostech a na mikroskopických vzdálenostech se dostáváme do oblasti \emph{kvantové mechaniky}, kde klasické pojmy přesné polohy a trajektorie částice ztrácejí svůj původní smysl.

\newpage
\pagenumbering{arabic}

% ---------------------------------------------------------------------------
% Hlavní kapitoly
% ---------------------------------------------------------------------------
\chapter{Popis polohy a směru v prostoru}

\section{Vektory}

velkost

aritmetika

skalarni a vektorovy soucin, geometricka interpretace

pravotocivy system, pravidlo prave ruky

\section{Systémy souřadnic}

konkretne numericke vyjadreni vektoru

\subsection{Kartezke souradnice}

\subsection{Cylindricke (polarni) souradnice}

\subsection{Sfericke souradnice}

\subsubsection{test}

\section{Transformace souradnic}

posuv

nasobeni vektoru matici - skalovani a rotace


\chapter{Kinematika a Newtonova dynamika}

% \section{Kinematika hmotného bodu}
% Poloha částice v trojrozměrném eukleidovském prostoru v čase $t$ je popsaná polohovým vektorem $\vect{r}(t) \in \mathbb{R}^3$. V ortonormální kartézské bázi $\{\vect{e}_x, \vect{e}_y, \vect{e}_z\}$ máme vyjádření:
% \begin{equation}
%     \vect{r}(t) = x(t)\vect{e}_x + y(t)\vect{e}_y + z(t)\vect{e}_z.
% \end{equation}

% Okamžitá rychlost $\vect{v}(t)$ a zrychlení $\vect{a}(t)$ jsou definovány pomocí časových derivací polohového vektoru:
% \begin{equation}
%     \vect{v}(t) \coloneqq \odv{\vect{r}}{t} = \dot{\vect{r}}(t), \qquad
%     \vect{a}(t) \coloneqq \odv{\vect{v}}{t} = \dodv{\vect{r}}{t} = \ddot{\vect{r}}(t).
% \end{equation}

% \begin{definice}[Hybnost a síla]
% Pro hmotný bod o konstantní hmotnosti $m$ pohybující se rychlostí $\vect{v}$ definujeme \textbf{hybnost}:
% \begin{equation}
%     \vect{p} \coloneqq m \vect{v}.
% \end{equation}
% Působí-li na bod celková síla $\vect{F}$, je časová změna hybnosti rovna této síle.
% \end{definice}

% \section{Newtonovy pohybové zákony}
% Základem newtonovské formulace jsou tři pohybové postuláty.

% \begin{veta}[Druhý Newtonův pohybový zákon]
% Časová změna hybnosti hmotného bodu je přímo úměrná výsledné působící síle $\vect{F}$ a má její směr:
% \begin{equation}
%     \odv{\vect{p}}{t} = \vect{F}(\vect{r}, \dot{\vect{r}}, t).
% \end{equation}
% V případě konstantní hmotnosti $m = \text{konst.}$ přechází rovnice do známého tvaru:
% \begin{equation}
%     m \ddot{\vect{r}} = \vect{F}.
% \end{equation}
% \end{veta}

% \begin{priklad}[Jednorozměrný lineární harmonický oscilátor]
% Uvažujme hmotný bod o hmotnosti $m$ na pružině s tuhostí $k > 0$ bez tlumení. Pohybová rovnice má tvar:
% \begin{equation}
%     m \ddot{x} + k x = 0 \iff \ddot{x} + \omega_0^2 x = 0, \quad \text{kde } \omega_0 = \sqrt{\frac{k}{m}}.
% \end{equation}
% Obecné řešení této lineární diferenciální rovnice druhého řádu je:
% \begin{equation}
%     x(t) = A \cos(\omega_0 t + \varphi_0),
% \end{equation}
% kde amplituda $A \ge 0$ a fázový posuv $\varphi_0 \in [0, 2\pi)$ plynou z počátečních podmínek $x(0)$ a $\dot{x}(0)$.
% \end{priklad}

% \begin{poznamka}
% Při transformaci do neinerciální vztažné soustavy (např. rotující úhlovou rychlostí $\vect{\omega}$) vystupují na pravé straně pohybové rovnice setrvačné síly: odstředivá síla $-m\vect{\omega}\crossprod(\vect{\omega}\crossprod\vect{r})$, Coriolisova síla $-2m(\vect{\omega}\crossprod\vect{v})$ a Eulerova síla $-m(\dot{\vect{\omega}}\crossprod\vect{r})$.
% \end{poznamka}

\chapter{Zákony zachování a pohyb v centrálním poli}

% \section{Integrály pohybu}
% Základní zákony zachování v mechanice bezprostředně souvisí se symetriemi časoprostoru (teorém Noetherové):
% \begin{itemize}
%     \item \textbf{Zachování mechanické energie}: plyne z homogenity času v konzervativním silovém poli $\vect{F} = -\nabla V(\vect{r})$:
%     \begin{equation}
%         E = T + V = \frac{1}{2}m\dot{\vect{r}}^2 + V(\vect{r}) = \text{konst.}
%     \end{equation}
%     \item \textbf{Zachování celkové hybnosti}: plyne z homogenity prostoru pro izolovanou soustavu částic.
%     \item \textbf{Zachování momentu hybnosti}: plyne z izotropie prostoru. Pro centrální pole $\vect{F} \parallel \vect{r}$ platí:
%     \begin{equation}
%         \vect{L} \coloneqq \vect{r} \crossprod \vect{p} = \text{konst.}
%     \end{equation}
% \end{itemize}

% \section{Pohyb v centrálním gravitačním poli}
% Vzhledem k zachování $\vect{L}$ probíhá pohyb v rovině kolmé na $\vect{L}$. V polárních souřadnicích $(r, \varphi)$ se úloha redukuje na jednorozměrný pohyb v efektivním potenciálu:
% \begin{equation}
%     E = \frac{1}{2}m\dot{r}^2 + V_{\text{eff}}(r), \qquad V_{\text{eff}}(r) = -\frac{\alpha}{r} + \frac{L^2}{2mr^2}.
% \end{equation}


% ---------------------------------------------------------------------------
% Literatura
% ---------------------------------------------------------------------------
\begin{thebibliography}{9}
\bibitem{landau}
    L.~D.~Landau, E.~M.~Lifshitz:
    \emph{Mechanics (Course of Theoretical Physics, Volume 1)}.
    Butterworth-Heinemann, 3. vydání, 1976.
\bibitem{goldstein}
    H.~Goldstein, C.~Poole, J.~Safko:
    \emph{Classical Mechanics}.
    Addison-Wesley, 3. vydání, 2001.
\bibitem{kvasnica}
    J.~Kvasnica:
    \emph{Mechanika}.
    Academia, Praha, 1988.
\end{thebibliography}


\end{document}
